\section{Introduction}

The rapid growth of digital media has produced a vast amount of structured information about movies, actors, directors, production companies, genres, and ratings. However, this information is often scattered across isolated databases and tabular datasets, making it difficult to query and connect semantically related entities. Knowledge graphs address this problem by representing entities and their relationships in an interconnected graph structure, supporting integration, exploration, and reasoning across heterogeneous information sources~\citep{hogan2021knowledge}. This perspective is consistent with the Linked Data principles of using dereferenceable identifiers and links to related resources~\citep{bernerslee2006linkeddata}.

This project develops a Linked Open Data (LOD) application for a movie knowledge graph, with the goal of transforming movie data into a reusable semantic resource. The system collects movie information from The Movie Database (TMDB), whose API provides access to movie, television, actor, image, language, and country data~\citep{themoviedatabase2025api}. The collected information is stored as CSV data and then transformed by this project into RDF triples. The project also defines a domain ontology based on the shared Schema.org vocabulary~\citep{schemaorg2025} and links shared entities to external knowledge bases such as DBpedia and Wikidata~\citep{dbpedia2025,wikidata2025}. RDF provides the graph-based data model used to express these resources as subject--predicate--object triples~\citep{w3c2014rdf11}. The resulting graph is intended to support semantic querying through SPARQL and to provide a structured, interoperable view of the movie domain.

The main objective of this work is to demonstrate a complete pipeline for building a movie knowledge graph in accordance with LOD principles. Our method follows five main stages: (1) collecting relevant movie data from TMDB, (2) transforming the collected CSV data into RDF in Turtle notation using stable URIs, (3) building an OWL ontology that defines the domain classes, properties, restrictions, and relationships, (4) establishing links to DBpedia and Wikidata to improve interoperability and achieve a higher level of Linked Open Data integration, and (5) providing a SPARQL interface through Apache Jena Fuseki for querying the resulting graph~\citep{w3c2013sparql}. An optional static web application is also provided to support browsing and exploration of the generated data.

The scope of the project is intentionally limited to the movie domain. It covers canonical entities such as movies, people, organizations, genres, countries, languages, and aggregate ratings, while preserving relationships such as cast membership, directorship, production company affiliation, and publication date. The project does not aim to cover all cultural or entertainment data globally; instead, it focuses on a representative and technically meaningful subset of the movie domain that can be modeled consistently using RDF, OWL, and Schema.org vocabulary~\citep{schemaorg2025}. This scope allows the system to remain reproducible, explainable, and aligned with the broader goals of semantic web applications.

This remainder of this report is organized as follows. ...